% answers_short.tex —— 2019 年·北京卷·理科 · 简短答案速查
% 由 tools/make_short_answers.py 自 answers.tex 的 \gksolhead/\gkanswer 提取，
% 并与 2019/2019_数学真题_北京卷理_答案要点.md 逐题比对后生成，勿手改（改源 MD 后重跑）。

\begin{gknotice}
\textbf{使用说明}：本节只给最终答案便于快速核对，完整推导见解析版。答案为编写者独立推导，\textbf{非官方标准答案}。
\end{gknotice}

\gkpart{选择题}
\begin{center}\begin{tabular}{ccccccccc}
\toprule
题号 & 1 & 2 & 3 & 4 & 5 & 6 & 7 & 8 \\
\midrule
答案 & D & B & D & B & C & A & C & C \\
\bottomrule
\end{tabular}\end{center}

\gkpart{填空题}
\begin{itemize}[leftmargin=2.6em]
  \item[\textbf{9.}] $\dfrac{\pi}{2}$
  \item[\textbf{10.}] $0$；$-10$
  \item[\textbf{11.}] $28$
  \item[\textbf{12.}] ②③ $\implies$ ①（或 ①③ $\implies$ ②）
  \item[\textbf{13.}] $-1$；$(-\infty, 0]$（或 $a \leqslant 0$）
  \item[\textbf{14.}] ① $130$；② $15$
\end{itemize}

\gkpart{解答题}
\begin{itemize}[leftmargin=2.6em]
  \item[\textbf{15.}] (1) $b = 7, c = 5$；(2) $\dfrac{4\sqrt{3}}{7}$
  \item[\textbf{16.}] (1) 证明见解析；(2) $\dfrac{\sqrt{3}}{3}$；(3) 直线 $AG$ 在平面 $AEF$ 内，理由见解析
  \item[\textbf{17.}] (1) $0.4$；(2) 分布列见解析，$E(X) = 1$；(3) 能认为有变化，理由见解析
  \item[\textbf{18.}] (1) 抛物线 $C$ 的方程为 $x^2 = -4y$，准线方程为 $y = 1$；(2) 证明见解析，定点坐标为 $(0, 1)$ 和 $(0, -3)$
  \item[\textbf{19.}] (1) $y = x$ 或 $y = x - \dfrac{64}{27}$；(2) 证明见解析；(3) $a = -3$
  \item[\textbf{20.}] (1) $1, 3, 5, 6$（或任一符合条件的长度为 4 的递增子列）；(2) 证明见解析；(3) $a_n = \begin{cases} n + 1, & n \text{ \text{为奇数}}, \\ n - 1, & n \text{ \text{为偶数}} \end{cases}$（或 $a_n = n + (-1)^{n-1}$）
\end{itemize}
